\documentclass[10pt,a4paper]{amsart}
\usepackage[margin=19mm]{geometry}
\usepackage{amsmath,amssymb,mathtools}
\usepackage[scr=rsfs,cal=euler]{mathalfa}
\usepackage{xfrac}
\usepackage[unicode,hidelinks,bookmarks=false]{hyperref}
\newcommand{\ie}{i.e.\ }
\newcommand{\cf}{cf.\ }
\newcommand{\resp}{respectively\ }
\newcommand{\Subsection}{\subsection}
\providecommand{\nobreakdash}{\nobreak\hbox{-}}
\setlength{\emergencystretch}{2em}
\multlinegap0pt
\usepackage{amssymb}
\usepackage{xcolor}
\usepackage{tikz}
\newtheorem{theorem}{Theorem}
\newcommand{\Esub}[2]{\mathbb{E}_{#2}\left[ #1 \right]}
\newcommand{\abs}[1]{\left| #1 \right|}
\newcommand{\bracket}[1]{\left[ #1 \right]}
\newcommand{\paren}[1]{\left( #1 \right)}
\title{The Thermodynamic Costs of Simple Linear Regression}
\author{Samuel H. D'Ambrosia$^\dagger$, Sultan M. Daniels$^\dagger$, Michael R. DeWeese, and Anant Sahai}
\date{Source: arXiv:2605.19195v1 (CC BY 4.0)}
\begin{document}
\maketitle

\noindent\emph{Source and licence.} The Thermodynamic Costs of Simple Linear Regression. Version arXiv:2605.19195v1 (CC BY 4.0), distributed by arXiv under the Creative Commons Attribution 4.0 International licence, http://creativecommons.org/licenses/by/4.0/, as recorded in the arXiv licence metadata for this exact version and shown on the abstract page https://arxiv.org/abs/2605.19195v1. Full licence notice: \texttt{licenses/arxiv-2605.19195-CC-BY-4.0.txt}.\par\smallskip

\noindent\textit{Editorial scope.} Counted unit: the article's theorem thm:steps (source0.tex lines 1406--1412). The article's complete original proof of it is delivered with it (source0.tex lines 1414--1499, one \texttt{proof} environment of 86 lines). No mathematics is written, repaired or re-derived: the statement and the proof are the article's own lines, in the article's own notation, and no retained line is substituted or re-flowed.  Retained uncounted as the source-local context the proof needs: lines 353--356; lines 1160--1160; lines 1394--1404.  Nothing was omitted from any retained span, so the package carries no omission record.  No retained line contains a citation, so the wrapper carries no bibliography and no bibliography entry.  No retained line contains a figure environment, an image-inclusion command or drawing code, so no image is delivered and the package declares no asset dependency.  Editorial formatting only, in addition: an AMS \texttt{amsart} preamble that restates the article's own declarations and macros used by the retained lines, namely the environments \texttt{theorem} on separate counters, together with the standard packages those lines need.  The retained environments print in this document as Theorem 1 in order of appearance, whereas the article prints the same items with the numbering of its own full document.  The complete native source of the article as distributed in the arXiv e-print for this exact version is bundled unchanged as \texttt{sources/200901/source0.tex}.
\begin{equation}
    \label{eq:true_bins}
    \Delta(x) = \frac{2^{e_{fp}(x)+1} - 2^{e_{fp}(x)}}{2^{p-1}}= 2^{e_{fp}(x) - (p-1)} \text{ if }x \in \mathbb{U}.
\end{equation}

\subsubsection*{Approximation 1 -- Relating discrete and differential entropy for non-uniform bins}\label{app:FPN_app1}

\begin{equation}
\label{eq:special_error}
\varepsilon_{0}(\mathbf X)
\triangleq
\sum_{j=1}^{d}
\int_{\mathcal B_{0}}
f_{X_j}(x)
\left|
\log\bracket{\frac{\Delta_s(x)}{\Delta_0}}
\right|dx.
\end{equation}

\begin{theorem}[Error bound for smoothing the bin-size function]\label{thm:steps}
Let $\mathbf{X}\sim f$ be a $d$-dimensional random vector with the probability density $f$ supported on $\mathbb U^d$, and let $X_j$ denote its $j$th component. Let $\tilde H(\mathbf X_{fp})$ be the approximation from App.~\ref{app:FPN_app1} using the true midpoint bin-size function $\Delta$, and let $\tilde H_{s,\mathbb U}(\mathbf X_{fp})$ be the corresponding approximation on $\mathbb U^d$ using $\Delta_s$. Then
\begin{equation}
\abs{\tilde H(\mathbf X_{fp}) - \tilde H_{s,\mathbb U}(\mathbf X_{fp})} \leq \frac{d}{2} + \varepsilon_{0}(\mathbf X),
\end{equation}
where $\varepsilon_{0}(\mathbf X)$ is finite and defined in Eq.~\eqref{eq:special_error}.
\end{theorem}

\begin{proof}

There are 3 types of bins for which we will bound the error: bins on the interior of an exponent block, bins on the outer edges, and bins on the boundary between exponent blocks (regions for which a single exponent value applies). The bound can be shown by considering the ratio between bin sizes, case by case.

\paragraph{Interior Bins}
For interior bins strictly within each exponent block $e$, by Eq.~\eqref{eq:true_bins},
\begin{equation}
    \label{eq:real_bins}
    \Delta(x_j) = 2^{e-(p-1)} \quad \text{for } 2^e \leq |x_j| < 2^{e+1},
\end{equation}
and within this range $\Delta_s(x_j) = \frac{|x_j|}{\sqrt{2}} \cdot 2^{1-p}$ satisfies
\begin{equation}
    \begin{aligned}
    &\frac{2^e}{\sqrt{2}}2^{1-p} \leq \Delta_s(x_j) \leq \frac{2^{e+1}}{\sqrt{2}}2^{1-p}\\
    \implies &\frac{1}{\sqrt{2}} \leq \frac{\Delta_s(x_j)}{\Delta(x_j)} < \sqrt{2}.
    \end{aligned}
\end{equation}

\paragraph{Outer clipping bins}

For the outer clipping bins, by symmetry it suffices to analyze $[\frac{u_K + u_{K-1}}{2}, 2^{e_{\max}+1}-2^{e_{\max}-p}]$ which is the bin on positive side of the real line. When $p \geq 2$, we know that $u_K = 2^{e_{max}+1}(1 - 2^{-p})$, and $u_{K-1} = 2^{e_{max}+1}(1 - 2^{1-p})$, so the outer clipping bin size is \begin{equation}
\Delta(x_j) = 2^{e_{max} + 1} - 2^{e_{max} - p}- \frac{u_K + u_{K-1}}{2} = 2^{e_{max}}\paren{2^{1-p}} = 2^{e_{max}-p + 1}
\end{equation}
Now, we have
\begin{equation}
\begin{aligned}
  &\frac{(u_K + u_{K-1})}{2\sqrt{2}}2^{1-p} \leq \Delta_s(x_j) \leq \frac{2^{e_{max} + 1} - 2^{e_{max} - p}}{\sqrt{2}}2^{1-p}\\
  \implies &\frac{ 2^{e_{max}}(1 - 2^{-1-p} - 2^{-p})}{2^{e_{max} - p}\sqrt{2}}2^{1-p} \leq \frac{\Delta_s(x_j)}{\Delta(x_j)} \leq \frac{2^{e_{max}}(1 - 2^{-1 - p})}{2^{e_{max} - p}\sqrt{2}}2^{1-p}\\
  \implies &\frac{ 2(1 - 2^{-1-p} - 2^{-p})}{\sqrt{2}} \leq \frac{\Delta_s(x_j)}{\Delta(x_j)} \leq \frac{2(1- 2^{-1-p})}{\sqrt{2}}\\
  \implies &\sqrt{2}(1 - 2^{-1-p} - 2^{-p}) \leq \frac{\Delta_s(x_j)}{\Delta(x_j)} \leq \sqrt{2}(1- 2^{-1-p}).
\end{aligned}
\end{equation}
When $p = 1$, $u_K = 2^{e_{max}}$ and $u_{K-1} = 2^{e_{max}-1}$. The true bin width is 
\begin{equation}
\Delta(x_j) = 2^{e_{max}+1} - 2^{e_{max}-1} - \frac{1}{2}(2^{e_{max}} + 2^{e_{max}-1}) = 3\cdot2^{e_{max}-2}.
\end{equation}
Using the same bounding technique, we have
\begin{equation}
\begin{aligned}
  &\frac{2^{e_{max}} + 2^{e_{max}-1}}{2\sqrt{2}} \leq \Delta_s(x_j) \leq \frac{2^{e_{max}+1} - 2^{e_{max}-1}}{\sqrt{2}}\\
  \implies &\frac{3\cdot 2^{e_{max}-2}}{3\cdot2^{e_{max}-2}\sqrt{2}} \leq \frac{\Delta_s(x_j)}{\Delta(x_j)} \leq \frac{3\cdot 2^{e_{max}-1}}{3\cdot2^{e_{max}-2}\sqrt{2}}\\
  \implies &\frac{ 1}{\sqrt{2}} \leq \frac{\Delta_s(x_j)}{\Delta(x_j)} \leq \frac{2}{\sqrt{2}}.
\end{aligned}
\end{equation}
For the outer clipping bins, we see that for all $p\geq 1$, $\frac{1}{\sqrt{2}} \leq \frac{\Delta_s(x_j)}{\Delta(x_j)} \leq \sqrt{2}$.

\paragraph{Exponent boundary bins}

At exponent boundaries, let the last representable value of the $e$-th exponent block be $u_A = 2^{e+1} - 2^{e-(p-1)}$. The first representable value of the $e+1$-th exponent block be $u_B = 2^{e+1}$, and the second representable value of the $e+1$-th exponent block be $u_C = 2^{e+1} + 2^{e+1 - (p-1)}$. The left boundary of the midpoint-quantization bin for $u_B$ is $m_L = \frac{u_A + u_B}{2} = 2^{e+1}(1 - 2^{-p-1})$ while the right boundary is $m_R = \frac{u_B + u_C}{2} = 2^{e+1}(1+2^{-p})$. When $x_j \in [m_L, m_R)$, the bin width is
\begin{equation}\label{eq:exp_boundary}
\Delta(x_j) = m_R - m_L = 3 \cdot 2^{e-p},
\end{equation}
and the ratio $\Delta_s(x_j)/\Delta(x_j)$ is bounded by
\begin{equation}
    \begin{aligned}
    &\frac{2^{e+1}(1 - 2^{-p-1})}{\sqrt{2}}2^{1-p} \leq \Delta_s(x_j) \leq \frac{2^{e+1}(1+2^{-p})}{\sqrt{2}}2^{1-p}\\
    \implies & \frac{2^{e+1}(1 - 2^{-p-1})}{3 \cdot 2^{e-p}\sqrt{2}}2^{1-p} \leq \frac{\Delta_s(x_j)}{\Delta(x_j)} \leq \frac{2^{e+1}(1+2^{-p})}{3 \cdot 2^{e-p}\sqrt{2}}2^{1-p}\\
    \implies & \frac{2\sqrt{2}(1 - 2^{-p-1})}{3} \leq \frac{\Delta_s(x_j)}{\Delta(x_j)}  \leq \frac{2\sqrt{2}(1+2^{-p})}{3}.
    \end{aligned}
\end{equation}
When $p = 1$ the lower bound is equal to $1/\sqrt{2}$ and the upper bound is equal to $\sqrt{2}$. The lower bound is monotonically increasing in $p$ while the upper bound is monotonically decreasing in $p$, and both converge to $2\sqrt{2}/3$ as $p \rightarrow \infty$, which is between $1/\sqrt{2}$ and $\sqrt{2}$. This means both bounds lie within $[1/\sqrt{2},\, \sqrt{2}]$, so for within exponent block bins and exponent boundary bins,
\begin{equation}
    -\frac{1}{2} \leq \log\left[\frac{\Delta_s(x_j)}{\Delta(x_j)}\right] < \frac{1}{2}.
\end{equation}

\paragraph{Bound on entropy difference from bin ratios}


Hence, for every $x_j\in \mathbb U\setminus \mathcal B_{0}$,
\begin{equation}
\abs{\log\bracket{\frac{\Delta_s(x_j)}{\Delta(x_j)}}} \leq \frac{1}{2}.
\end{equation}

The remaining points are exactly the special bins collected in $\mathcal B_{0}$, where we do not claim a uniform pointwise bound and instead keep their contribution explicitly:
\begin{equation}
\begin{aligned}
&\left|
\tilde H(\mathbf X_{fp}) - \tilde H_{s,\mathbb U}(\mathbf X_{fp}) \right| = \Bigg | -\int_{\mathbb U^d} f_{\mathbf X}(\mathbf x)\,\log\left[f_{\mathbf X}(\mathbf x)\prod_{j=1}^d \Delta(x_j)\right]d\mathbf x  \\ & \qquad \qquad \qquad \qquad \qquad
+ \int_{\mathbb U^d} f_{\mathbf X}(\mathbf x)\,\log\left[f_{\mathbf X}(\mathbf x)\prod_{j=1}^d \Delta_s(x_j)\right]d\mathbf x \Bigg |\\
&\left|
\tilde H(\mathbf X_{fp}) - \tilde H_{s,\mathbb U}(\mathbf X_{fp}) \right| = \abs{\sum_{j=1}^{d} \Esub{\log\bracket{\frac{\Delta_s(x)}{\Delta(x)}}}{X_j}} \\
&\leq\sum_{j=1}^{d} \int_{\mathbb U\setminus \mathcal B_{0}} f_{X_j}(x) \abs{\log\bracket{\frac{\Delta_s(x)}{\Delta(x)}}}dx + \varepsilon_{0}(\mathbf X) \leq\frac d2 + \varepsilon_{0}(\mathbf X).
\end{aligned}
\end{equation}
The point $x=0$ has measure zero, and $\varepsilon_{0}(\mathbf X)$ is finite since each marginal density is continuous on $\mathbb U$ while $|\log|x||$ is locally integrable near $0$.
\end{proof}

\end{document}
